\chapter{Infinite-mode states and local stabilization}
\label{ch:vol2-infinite-mode-states}

A current commutator does not determine a state. For one oscillator,
\([a,a^*]=1\), the vacuum has expectation zero on \(a^*a\).
A thermal state at inverse temperature \(T\) has expectation
\((e^T-1)^{-1}\) on the same operator. Both give expectation one
on the commutator. The singular current coefficient can therefore
survive a change of preparation while its regular correlations change.

Chapter~\ref{ch:vol2-gaussian-elimination} constructed finite Gaussian
states and their higher operations. Passing to infinitely many modes
requires a domain, an energy estimate and a normalization. We construct
these data on a Heisenberg Fock space. A local cubic interaction then
separates two possible conclusions: every perturbative coefficient
exists and sews, although the numerical cubic heat operator is unbounded.
A local quartic correction gives positive states. Its coefficient
determines whether those states approach the free state continuously.

\section{The common Hilbert space}

Fix a numerical level \(\lambda>0\). Positivity in this chapter
concerns complex Hilbert spaces with this fixed real parameter.
It is separate from an identity over a formal quantum coefficient ring.
Let
\[
 \mathcal H_1=\ell^2(\mathbb N_{>0}),\qquad
 \mathcal F=\Gamma_s(\mathcal H_1)
 =\widehat{\bigoplus}_{r\geq0}\operatorname{Sym}^r\mathcal H_1.
\]
The inner product is conjugate linear in its first argument.
The occupation vectors \(\lvert\mathbf m\rangle\), indexed by
finite-support sequences \(\mathbf m=(m_1,m_2,\ldots)\), form an
orthonormal basis. Define
\[
 a_n\lvert\mathbf m\rangle
 =\sqrt{m_n}\lvert\mathbf m-\mathbf e_n\rangle,\qquad
 a_n^*\lvert\mathbf m\rangle
 =\sqrt{m_n+1}\lvert\mathbf m+\mathbf e_n\rangle.
\]
Thus \([a_m,a_n^*]=\delta_{mn}\). Write
\[
 \alpha_n=\sqrt n\,a_n,\qquad
 \alpha_{-n}=\sqrt n\,a_n^*,\qquad
 J_n=\sqrt\lambda\,\alpha_n\quad(n\ne0),
 \qquad\alpha_0=J_0=0.
\]
The current relations are
\[
 [\alpha_m,\alpha_n]=m\delta_{m+n,0},\qquad
 J_n^*=J_{-n},\qquad
 [J_m,J_n]=\lambda m\delta_{m+n,0}.
\]

The vacuum \(\Omega=\lvert\mathbf0\rangle\) has energy zero.
The energy operator and particle-number operator are
\[
 H_0=\sum_{n\geq1}n a_n^*a_n
     =\sum_{n\geq1}\alpha_{-n}\alpha_n,
 \qquad N=\sum_{n\geq1}a_n^*a_n\leq H_0.
\]
These are diagonal self-adjoint operators on their maximal domains.
For example,
\[
 \operatorname{Dom}(H_0)
 =\left\{\sum_{\mathbf m}v_{\mathbf m}\lvert\mathbf m\rangle:
   \sum_{\mathbf m}\left(\sum_n n m_n\right)^2
              |v_{\mathbf m}|^2<\infty\right\}.
\]
Let \(\mathcal F_E\) be the energy-\(E\) eigenspace and \(P_E\)
its orthogonal projection. Its dimension is the partition number
\(p(E)\). The finite-energy space
\(\mathcal F_{\mathrm{fin}}=\bigoplus_{E\geq0}\mathcal F_E\)
is a common core for the finite-mode operators. Their closed extensions
are used whenever an infinite vector is an input.

For \(f\in\mathcal H_1\), the smeared oscillators are
\[
 a(f)=\sum_n\overline{f_n}a_n,\qquad
 a^*(f)=\sum_n f_n a_n^*.
\]
They act on the algebraic direct sum of the particle-number spaces,
which permits arbitrary square-summable one-particle vectors.
On that domain,
\[
 \|a(f)v\|\leq\|f\|_2\|N^{1/2}v\|,
 \qquad
 \|a^*(f)v\|\leq\|f\|_2\|(N+1)^{1/2}v\|.
\]
Successive application bounds a polynomial of degree \(d\) in finitely
many such oscillators by a constant times \((N+1)^{d/2}\).
This estimate specifies the unbounded observables whose thermal
expectations will be used below.

\section{A prepared trace and its current kernel}

Choose a cylinder length \(T>0\) and put \(q=e^{-T}\).
Its evolution operator is \(K_T=e^{-TH_0}\). This choice supplies
the preparation before any comparison with a bulk theory.

\begin{theorem}\label{thm:vol2-fock-thermal-state}
The operator \(K_T\) is positive and trace class. If
\(E_1 e_n=n e_n\) on \(\mathcal H_1\), then
\begin{equation}\label{eq:vol2-fock-partition}
 Z_0(T)=\operatorname{Tr}_{\mathcal F}K_T
 =\prod_{n\geq1}(1-e^{-Tn})^{-1}
 =\det_{\mathcal H_1}(1-e^{-TE_1})^{-1}.
\end{equation}
The density \(\rho_T=K_T/Z_0(T)\) defines a positive normalized
state on the bounded operators on \(\mathcal F\). Every polynomial
in finitely many square-summably smeared oscillators also has a
thermal expectation
\[
 \varphi_T(A)=Z_0(T)^{-1}\operatorname{Tr}(A K_T).
\]
For the radially ordered current \(J(z)=\sum_nJ_nz^{-n-1}\),
\begin{equation}\label{eq:vol2-thermal-current}
 \varphi_T(J(z)J(w))
 =\frac{\lambda}{zw}\sum_{n\geq1}
 \frac{n\bigl((w/z)^n+q^n(z/w)^n\bigr)}{1-q^n},
 \qquad q<|w/z|<1.
\end{equation}
The series and every fixed holomorphic derivative converge normally
on compact subannuli. Near \(z=w\ne0\), its singular part is
\(\lambda/(z-w)^2\).
\end{theorem}

\begin{proof}
At a mode cutoff \(M\), independent occupation sums give
\(Z_{0,M}(T)=\prod_{n=1}^M(1-q^n)^{-1}\). Moreover,
\begin{equation}\label{eq:vol2-partition-tail}
 0\leq\log\frac{Z_0(T)}{Z_{0,M}(T)}
 \leq\frac{q^{M+1}}{(1-q)(1-q^{M+1})}.
\end{equation}
Indeed \(-\log(1-x)\leq x/(1-x)\), and summing the resulting
geometric bound proves convergence of the logarithmic product.
Monotone summation over all occupation vectors gives its equality
with the trace. The eigenvalues \(q^n\) of \(e^{-TE_1}\) are
summable, so its Fredholm determinant is this same convergent product.

For every \(r\geq0\), the operator \((H_0+1)^r e^{-TH_0}\)
is trace class: the polynomial factor is bounded by a constant
times \(e^{TE/2}\) on energy \(E\), and the remaining trace is
\(Z_0(T/2)\). The oscillator bounds therefore make \(A K_T\)
trace class. They also make \(A K_{T/2}\) Hilbert--Schmidt.
Consequently the expectation of \(A^*A\) is its squared
Hilbert--Schmidt norm divided by \(Z_0(T)\), and is nonnegative.
This extends the positive state to the stated polynomial domain.

The one-mode geometric sums give
\[
 \varphi_T(J_nJ_{-n})=\frac{\lambda n}{1-q^n},\qquad
 \varphi_T(J_{-n}J_n)=\frac{\lambda n q^n}{1-q^n}
 \quad(n>0).
\]
They imply \eqref{eq:vol2-thermal-current}. On a compact subannulus,
both \(|w/z|\) and \(q|z/w|\) are bounded strictly below one.
The denominator is bounded below by \(1-q\), and derivatives add
only polynomial factors in \(n\). Split the first numerator into
its vacuum contribution and its \(q^n\)-weighted remainder.
The vacuum sum \(\sum_{n\geq1}n(w/z)^n\) gives
\(\lambda/(z-w)^2\); the two remaining sums extend across \(z=w\).
\end{proof}

At a cutoff, the normalized-oscillator contractions are bounded by
\((1-q)^{-1}\|f\|_2\|h\|_2\). Wick expansion therefore gives
the same polynomial expectations by a mode limit. The estimates
prove convergence of an actual density and determinant, as well as
convergence of its correlations.

A fixed charge \(b\in\mathbb R\) can be retained by setting
\(J_0=b\) and adding \(b^2/(2\lambda)\) to the energy.
The partition function becomes \(e^{-Tb^2/(2\lambda)}Z_0(T)\),
the mean current is \(b/z\), and the connected two-point function
is unchanged. A sum over charges additionally requires a measure
or lattice and a convergence domain.

\section{Sewing before normalization}

For \(z\in\mathcal H_1\), define the exponential vector
\[
 \mathrm e(z)=\bigoplus_{r\geq0}\frac{z^{\otimes r}}{\sqrt{r!}}.
\]
Its norm squared is \(e^{\|z\|_2^2}\). The cylinder matrix element is
\begin{equation}\label{eq:vol2-coherent-cylinder}
 K_T(z,\overline w)
 =\langle\mathrm e(w),K_T\mathrm e(z)\rangle
 =\exp\!\left(\sum_{n\geq1}e^{-Tn}z_n\overline{w_n}\right).
\end{equation}
Cauchy--Schwarz makes the exponent convergent.

\begin{proposition}\label{prop:vol2-fock-sewing}
For \(T_1,T_2>0\), coherent cylinder kernels sew by
\(K_{T_1}\circ K_{T_2}=K_{T_1+T_2}\). Closing the composite by
the trace gives \(Z_0(T_1+T_2)\).
\end{proposition}

\begin{proof}
For a bounded one-particle contraction \(B\), second quantization is
\(\Gamma(B)=\bigoplus_{r\geq0}\operatorname{Sym}^rB\). It gives
\[
 \Gamma(e^{-T_1E_1})\Gamma(e^{-T_2E_1})
 =\Gamma(e^{-(T_1+T_2)E_1}).
\]
The finite-dimensional coherent Gaussian integral computes this
identity at each mode cutoff. Its limit is Fock-space composition.
The high-mode vacuum projections commute with the heat operators,
and the omitted positive trace tends to zero by
\eqref{eq:vol2-partition-tail}. Thus taking the trace of the
composite gives the stated determinant. Associativity is already
the associativity of operator composition.
\end{proof}

The infinite-mode pairing here is a limit of cylindrical Wick
pairings. No translation-invariant measure on \(\ell^2\) is used.
Normalization is performed after the required trace closure:
\[
 \rho_{T_1}\rho_{T_2}
 =\frac{Z_0(T_1+T_2)}{Z_0(T_1)Z_0(T_2)}\rho_{T_1+T_2}.
\]
The scalar factor is part of sewing data.

The thermal state also satisfies, for finite-mode oscillator
polynomials \(A,B\),
\begin{equation}\label{eq:vol2-fock-kms}
 \varphi_T(AB)
 =\varphi_T\bigl(B e^{-TH_0}A e^{TH_0}\bigr).
\end{equation}
This follows by cyclicity of the trace on the common domain, or
directly from the occupation sums. The conjugation multiplies
\(\alpha_n\) by \(e^{Tn}\). For general infinite smears it
therefore requires a further domain condition.
Equation~\eqref{eq:vol2-fock-kms} is the thermal, or KMS, identity.
Ordinary traciality fails already on \(a_na_n^*\) and \(a_n^*a_n\).
A cyclic chain class from Chapter~\ref{chap:vol2-cyclic-traces-wheels}
and this scalar state are consequently different constructions.

\section{The stress tensor and the cylinder constant}

The vacuum has contraction
\(\langle\Omega,J(z)J(w)\Omega\rangle=\lambda/(z-w)^2\)
in \(|w|<|z|\). Subtract this contraction to define
\[
 T(z)=\frac1{2\lambda}{:}J(z)^2{:}
     =\frac12{:}\alpha(z)^2{:}
     =\sum_{m\in\mathbb Z}L_mz^{-m-2},
 \qquad
 L_m=\frac12\sum_{j\in\mathbb Z}{:}\alpha_j\alpha_{m-j}{:}.
\]
Normal ordering puts negative-index creators to the left.
Terms of index zero vanish. Each sum is finite on a finite-energy
vector, and \(L_0=H_0\).

\begin{proposition}\label{prop:vol2-current-virasoro}
On \(\mathcal F_{\mathrm{fin}}\),
\[
 [L_m,\alpha_n]=-n\alpha_{m+n},\qquad
 [L_m,L_n]=(m-n)L_{m+n}
       +\frac{m^3-m}{12}\delta_{m+n,0}.
\]
If \(z=f(u)\) is a local holomorphic coordinate change with
\(f'\ne0\), the stress tensor defined by subtraction in the
\(u\)-coordinate is
\begin{equation}\label{eq:vol2-stress-schwarzian}
 T_u(u)=f'(u)^2T_z(f(u))+\frac1{12}\{f,u\},
 \qquad
 \{f,u\}=\frac{f'''}{f'}-\frac32\left(\frac{f''}{f'}\right)^2.
\end{equation}
In particular \(f(u)=e^u\) gives cylinder Hamiltonian
\(H_0-1/24\).
\end{proposition}

\begin{proof}
Commuting a current through the two factors of \(L_m\) gives
the first identity. The Jacobi identity then shows that
\([L_m,L_n]-(m-n)L_{m+n}\) commutes with every current.
An operator with this property is scalar: it preserves the common
kernel of the positive modes, which is \(\mathbb C\Omega\), and
then has the same scalar action on all created vectors.
Its energy degree forces that scalar to vanish unless \(m+n=0\).
For \(m>0\), its vacuum expectation is
\[
 \|L_{-m}\Omega\|^2
 =\frac12\sum_{j=1}^{m-1}j(m-j)
 =\frac{m^3-m}{12}.
\]
The two contractions retain the factor \(1/2\), including when
\(j=m-j\). Antisymmetry gives the other signs of \(m\).

The current transforms as a one-form, so its pulled-back vacuum
contraction is \(\lambda f'(u)f'(v)/(f(u)-f(v))^2\).
Taylor expansion at \(v=u\) gives
\[
 \frac{f'(u)f'(v)}{(f(u)-f(v))^2}
 =\frac1{(u-v)^2}+\frac16\{f,u\}+O(v-u).
\]
For completeness, writing \(a=f''/f'\), \(b=f'''/f'\), the
constant term is \(b/2-a^2+3a^2/4-b/3=b/6-a^2/4\).
Subtracting the \(u\)-coordinate pole and dividing by \(2\lambda\)
proves \eqref{eq:vol2-stress-schwarzian}.
Since \(\{e^u,u\}=-1/2\), its constant is \(-1/24\).
\end{proof}

The conformal cylinder convention consequently has partition function
\[
 \operatorname{Tr}e^{-T(H_0-1/24)}=e^{T/24}Z_0(T).
\]
The scalar factor respects sewing and cancels from normalized
expectations. It arises from the specified stress subtraction.
Identifying it with a determinant in another field theory requires
that theory's regulator, finite term, boundary condition and sewing
normalization to be compared.

\section{Formal interactions on two coefficient domains}

Let \(V\) be a finite linear combination of normally ordered
polynomials in \(a(f_j),a^*(f_j)\), for finitely many
\(f_j\in\ell^2\). Define a formal coupling series by
\begin{equation}\label{eq:vol2-state-dyson}
 \begin{split}
 U_T(g)&=\sum_{k\geq0}g^k U_T^{(k)},\qquad U_T^{(0)}=K_T,\\
 U_T^{(k)}&=(-1)^k\int_{0\leq t_1\leq\cdots\leq t_k\leq T}
 e^{-(T-t_k)H_0}V e^{-(t_k-t_{k-1})H_0}\cdots\\
 &\hspace{24mm}\cdots V e^{-t_1H_0}\,dt_1\cdots dt_k.
 \end{split}
\end{equation}
For this coefficient domain, the definition first uses Wick matrix
elements and traces at a mode cutoff, followed by their limits.

\begin{proposition}\label{prop:vol2-smeared-dyson}
Every thermal trace and coherent-vector matrix element of each
coefficient in \eqref{eq:vol2-state-dyson} has a mode-cutoff limit.
The same holds with a fixed polynomial smeared insertion.
In these coefficientwise sewn matrix elements,
\begin{equation}\label{eq:vol2-formal-state-sewing}
 U_{T_1}(g)U_{T_2}(g)=U_{T_1+T_2}(g).
\end{equation}
Dividing observable traces by \(\operatorname{Tr}U_T(g)\) defines
a normalized formal functional.
\end{proposition}

\begin{proof}
A time-ordered thermal contraction of normalized modes contains
one of
\[
 \frac{e^{-n\tau}}{1-e^{-nT}},\qquad
 \frac{e^{-n(T-\tau)}}{1-e^{-nT}},\qquad0\leq\tau\leq T.
\]
Both are bounded by \(C_T=(1-e^{-T})^{-1}\). A smeared contraction
is therefore at most \(C_T\|f\|_2\|h\|_2\), uniformly in the
mode cutoff and time. At fixed order there are finitely many monomial
choices and Wick pairings. Their bounds, multiplied by the simplex
volume \(T^k/k!\), are integrable. The contraction tails converge
by Cauchy--Schwarz, so dominated convergence gives the mode limits.

For coherent matrix elements, each external contraction contains
a one-particle heat contraction \(e^{-sE_1}\), with \(s\geq0\).
The free coherent factor is bounded by \(e^{\|z\|_2\|w\|_2}\).
These bounds prove the same assertion with external vectors.
A fixed insertion only adds finitely many contractions.

At a cutoff, split the time simplex at \(T_2\). Its pieces are
indexed by the number of insertions before the cut and give the
coefficient of \(g^k\) in \eqref{eq:vol2-formal-state-sewing}.
The bounds pass this identity to the limit. The denominator has
constant term \(Z_0(T)>0\), and hence a formal inverse.
\end{proof}

This domain permits infinitely many oscillator modes in a smear,
but only finitely many input directions. It does not permit a
point-current smear: its regularized one-particle norm contains
\[
 \sum_{n\geq1}n e^{-2\varepsilon n}
 =\frac{e^{-2\varepsilon}}{(1-e^{-2\varepsilon})^2}
 \sim\frac1{4\varepsilon^2}.
\]
Local spatial Wick polynomials require a different estimate.

For fixed \(n\in\mathbb Z\), \(d\geq1\) and polynomial \(P\),
define the local mode
\begin{equation}\label{eq:vol2-local-wick-mode}
 W_n=\sum_{n_1+\cdots+n_d=n}
 P(n_1,\ldots,n_d){:}\alpha_{n_1}\cdots\alpha_{n_d}{:}.
\end{equation}
Zero indices give zero terms. The polynomial \(P\) includes factors
from angular derivatives. Finite Fourier combinations of these
operators, together with the identity, form the second domain.

\begin{lemma}\label{lem:vol2-local-mode-bound}
The sum \eqref{eq:vol2-local-wick-mode} is finite on \(\mathcal F_E\),
maps it to \(\mathcal F_{E-n}\), and satisfies
\[
 \|W_nP_E\|\leq C_{n,d,P}(1+E+|n|)^{3d/2-1+\deg P}.
\]
The same estimate holds after discarding any set of its summands.
\end{lemma}

\begin{proof}
The annihilators act first. Their total index is at most \(E\)
whenever the monomial is nonzero. The creators' total absolute index
is that sum minus \(n\), hence at most \(E+|n|\).
Each contributing index is bounded by \(E+|n|\), and the sum
constraint determines the last from the first \(d-1\).
There are at most \(C_d(1+E+|n|)^{d-1}\) tuples.

Intermediate energies first decrease and then increase to \(E-n\).
The norm of each oscillator factor is at most
\((1+E+|n|)^{1/2}\): its square is bounded by the energy before
annihilation or after creation. Multiply the \(d\) bounds and
the polynomial bound for \(P\), then sum over tuples.
The exponent is \((d-1)+d/2+\deg P\). Since this argument bounds
each summand separately, it also bounds every partial sum.
\end{proof}

For example \(W_0=2H_0\) when \(d=2\) and \(P=1\).
Its one-particle compression has eigenvalues \(2,4,6,\ldots\).
A polynomial in finitely many square-summable smears instead has
one-particle compression equal to a scalar plus a bounded finite-rank
operator. The two domains therefore do not coincide.

\begin{theorem}\label{thm:vol2-local-mode-dyson}
Let \(V\) be a finite Fourier combination of local Wick modes.
For every \(T>0\) and \(k\geq0\), the coefficient
\(U_T^{(k)}\) has a unique trace-class extension from finite-energy
vectors. Mode-cutoff coefficients, embedded by the high-mode vacuum,
converge in trace norm. A fixed finite Fourier Wick-polynomial
insertion \(A\) also makes \(A U_T^{(k)}\) trace class.
The formal operators obey \eqref{eq:vol2-formal-state-sewing}, and
their normalized observable traces are well defined.
\end{theorem}

\begin{proof}
There are \(C,M,R\) such that
\[
 V\mathcal F_E\subset\bigoplus_{|r|\leq R}\mathcal F_{E+r},
 \qquad\|VP_E\|\leq C(1+E)^M.
\]
At order \(k\), every intermediate energy lies between
\(E-kR\) and \(E+kR\). For a fixed sequence of energy projections,
the heat factors are bounded by \(e^{-TE+TkR}\), uniformly over
the closed time simplex. The operator factors contribute at most
\(C^k(1+E+kR)^{Mk}\). Absorb the finite number of energy-shift
choices into \(C^k\). A block with source \(\mathcal F_E\)
has rank at most \(p(E)\), so
\begin{equation}\label{eq:vol2-local-dyson-tail}
 \|U_T^{(k)}\|_1
 \leq\frac{T^k}{k!}C^k e^{TkR}
        \sum_{E\geq0}p(E)e^{-TE}(1+E+kR)^{Mk}<\infty.
\end{equation}
The sum is finite because a polynomial energy weight is dominated
by \(e^{TE/2}\). Finite block sums are trace-norm continuous in
the simplex variables; the uniform summable tail gives a trace-norm
Bochner integral, including coincident-time faces.

For fixed \(E,k\), all contributing oscillator indices occur
below \(E+kR\), after enlarging this bound by the fixed finite
shift if necessary. Those blocks stabilize at a mode cutoff.
Lemma~\ref{lem:vol2-local-mode-bound} gives the same tail bound at
every cutoff, proving trace-norm convergence. An insertion only adds
fixed shifts and a polynomial factor to this estimate.

Subdivision of the simplex proves sewing on finite blocks.
The trace-norm convergence passes the identity to the full space,
using \(\|BC\|_1\leq\|B\|_1\|C\|\) and boundedness of each
trace-class coefficient. The free partition function is the
nonzero constant denominator. Uniqueness follows from the density
of finite-energy vectors.
\end{proof}

Both constructions are coefficientwise in \(g\). Neither estimate
asserts convergence as the coupling order tends to infinity.
The local-mode theorem concerns spatially local boundary operators
on the stated Fock carrier. A local interacting half-space theory
also requires its supported counterterms and quantum master equation.

\section{The cubic mode and its spectral obstruction}

Consider the self-adjoint finite-energy matrices of the cubic mode
\begin{equation}\label{eq:vol2-cubic-mode}
 \begin{split}
 V&=\frac16\operatorname{Res}_{z=0}z^2{:}\alpha(z)^3{:}\,dz\\
  &=\frac12\sum_{m,n\geq1}
   \bigl(\alpha_{-m}\alpha_{-n}\alpha_{m+n}
          +\alpha_{-(m+n)}\alpha_m\alpha_n\bigr).
 \end{split}
\end{equation}
It commutes with \(H_0\). The factor \(1/2\) comes from the
three placements of the uniquely signed mode in the residue.
In the polynomial realization \(\alpha_{-n}=p_n\),
\(\alpha_n=n\partial_{p_n}\), it becomes
\begin{equation}\label{eq:vol2-cut-and-join}
 V=\frac12\sum_{m,n\geq1}
 \left((m+n)p_mp_n\partial_{p_{m+n}}
       +mn p_{m+n}\partial_{p_m}\partial_{p_n}\right).
\end{equation}
This is the cut-and-join operator, with the normalization of
Mironov--Morozov--Natanzon \cite{MironovMorozovNatanzon2009}*{equation (51)}.

\begin{lemma}\label{lem:vol2-cubic-norm}
On energy \(E\), the restriction \(V_E\) has norm
\(E(E-1)/2\), and both extreme signs occur.
\end{lemma}

\begin{proof}
For a partition \(\mu=(1^{m_1}2^{m_2}\cdots)\) of \(E\), put
\(p_\mu=\prod_n p_n^{m_n}\) and
\(z_\mu=\prod_n n^{m_n}m_n!\). The Fock pairing is
\(\langle p_\mu,p_\nu\rangle=\delta_{\mu\nu}z_\mu\).
The characteristic map on class functions of \(S_E\),
\[
 \operatorname{ch}(f)=\sum_{\mu\vdash E}\frac{f(\mu)}{z_\mu}p_\mu,
\]
is an isometry for normalized counting measure. Convolution by the
sum of all transpositions preserves class functions. A transposition
joins two cycles or splits one. Counting its marked entries gives
the coefficients \(mn\) and \(m+n\) in
\eqref{eq:vol2-cut-and-join}, with the factor \(1/2\) for the two
orders. Thus this sum corresponds to \(V_E\).

Every transposition acts unitarily and self-adjointly by translation
on functions. The sum has norm at most \(\binom E2\).
The trivial and sign characters attain the positive and negative
values. Their characteristic images are the complete and elementary
symmetric polynomials \(h_E,e_E\), defined by
\[
 \sum_{E\geq0}h_Et^E
 =\exp\!\left(\sum_{n\geq1}\frac{p_nt^n}{n}\right),\qquad
 \sum_{E\geq0}e_Et^E
 =\exp\!\left(\sum_{n\geq1}\frac{(-1)^{n-1}p_nt^n}{n}\right).
\]
They give \(Vh_E=\binom E2h_E\) and
\(Ve_E=-\binom E2e_E\), proving equality in the norm bound.
\end{proof}

For nonzero real \(g\), one of these families has energy
\[
 (H_0+gV)\big|_{\mathrm{extreme},E}
 =E-\frac{|g|}{2}E(E-1)\longrightarrow-\infty.
\]
The self-adjoint direct sum \(H_0+gV\) is therefore unbounded
below, and \(e^{-T(H_0+gV)}\) is unbounded for every \(T>0\).
It cannot prepare a trace-class Gibbs state. The formal coefficients
from Theorem~\ref{thm:vol2-local-mode-dyson} nevertheless exist.
Since \([H_0,V]=0\), they are
\(U_T^{(k)}=(-T)^k e^{-TH_0}V^k/k!\).

The first response already sees the interaction. With
\(\langle A\rangle_q=\varphi_T(A)\), one has
\begin{equation}\label{eq:vol2-cubic-thermal-moment}
 \langle V^2\rangle_q
 =\sum_{m,n\geq1}
 \frac{mn(m+n)q^{m+n}}
      {(1-q^m)(1-q^n)(1-q^{m+n})}.
\end{equation}
To prove the formula, split \(V\) into its joining and splitting
terms. Only a join followed by a split, or the reverse order,
has complete thermal contractions. The two orders and the two
pairings of the \(m,n\) legs cancel the factor \(1/4\).
The same multiplicity remains when \(m=n\). Substitution of the
one-mode expectations proves the formula. Its convergence follows
from the exponential numerator and denominators bounded below by
\(1-q\). Its first coefficients are
\[
 2q^2+16q^3+60q^4+160q^5+360q^6+O(q^7).
\]
Oscillator parity sends \(V\) to \(-V\) and preserves the free
density, so all odd moments vanish. Hence the formal partition
function obeys
\[
 \log\frac{\operatorname{Tr}U_T(g)}{Z_0(T)}
 =\frac{g^2T^2}{2}\langle V^2\rangle_q+O(g^4).
\]

\section{Stabilization by the total energy}

A square of the total energy provides a first positive preparation.
For real \(\varepsilon,g\), define the self-adjoint block operator
\[
 H_{\varepsilon,g}=H_0+\varepsilon H_0^2+gV.
\]
The term \(H_0^2\) is a prescribed nonlocal interaction on the
circle. It changes the Hamiltonian whose state is being constructed.

\begin{theorem}\label{thm:vol2-nonlocal-stability}
For \(T>0\), the operator \(e^{-TH_{\varepsilon,g}}\) is positive
and trace class exactly when
\begin{equation}\label{eq:vol2-nonlocal-threshold}
 \varepsilon\geq |g|/2.
\end{equation}
On this region it gives a positive normalized state on bounded
operators, and its unnormalized heat kernels sew by addition of
cylinder lengths. For fixed \(T>0\), the partition function is
holomorphic in the complex tube
\(\operatorname{Re}\varepsilon>|\operatorname{Re}g|/2\).
\end{theorem}

\begin{proof}
Lemma~\ref{lem:vol2-cubic-norm} gives, on energy \(E\),
\begin{equation}\label{eq:vol2-nonlocal-energy}
 H_{\varepsilon,g}\geq
 (\varepsilon-|g|/2)E^2+(1+|g|/2)E.
\end{equation}
At equality in \eqref{eq:vol2-nonlocal-threshold}, this is still
a positive linear lower bound. Summation against \(p(E)\) proves
trace class by \eqref{eq:vol2-fock-partition}. Below the threshold,
choose the extreme eigenvector with the sign opposite to \(g\).
If \(g=0\), either extreme vector suffices. Its energy equals the
right side of \eqref{eq:vol2-nonlocal-energy} and tends to
\(-\infty\), so the heat operator is unbounded.

Spectral calculus gives positivity and sewing for the allowed
parameters. In a common eigenbasis of \(H_0,V\), each exponential
eigenvalue is entire in \((\varepsilon,g)\). On a compact subset
of the tube, its absolute value is at most
\(e^{-cE^2-TE}\) for some \(c>0\). Every parameter derivative
adds only a polynomial energy factor. Uniform summability proves
holomorphy of the trace.
\end{proof}

This convergence has an explicit tail. Put
\(c=\varepsilon-|g|/2\geq0\). For an integer \(M\geq0\),
\begin{equation}\label{eq:vol2-stabilized-tail}
 \sum_{E>M}\operatorname{Tr}_{\mathcal F_E}e^{-TH_{\varepsilon,g}}
 \leq Z_0(T/2)
 \frac{e^{-T(M+1)/2-Tc(M+1)^2}}
      {1-e^{-T/2-Tc(2M+3)}}.
\end{equation}
Indeed \(p(E)\leq Z_0(T/2)e^{TE/2}\). The remaining dominating
series has first term as displayed and a decreasing ratio, bounded
by its first ratio. The estimate includes the equality case \(c=0\).
At fixed nonzero real \(g\), the stabilizer cannot be removed
within this family while retaining trace class.

\section{The unchanged cubic series and its Borel boundary}

The original cubic formal trace has coefficients
\begin{equation}\label{eq:vol2-cubic-formal-trace}
 Z(T,g)=\sum_{n\geq0}a_n(T)g^n,
 \qquad
 a_n(T)=\frac{(-T)^n}{n!}\operatorname{Tr}(e^{-TH_0}V^n).
\end{equation}
Its odd coefficients vanish and its even coefficients are positive.
Define its ordinary Borel transform by
\[
 \widehat Z_T(\zeta)=\sum_{n\geq0}\frac{a_n(T)}{n!}\zeta^n.
\]
The Borel variable \(\zeta\) is a complex scalar independent of
the current coordinate and of the coupling.

\begin{theorem}\label{thm:vol2-cubic-borel-radius}
For every \(T>0\), \eqref{eq:vol2-cubic-formal-trace} has radius
zero in \(g\). Its Borel transform has exact radius \(T/2\)
and diverges at both real endpoints \(\zeta=\pm T/2\).
\end{theorem}

\begin{proof}
Fix \(0<\nu<T\) and write \(t=T-\nu\). Since
\(p(E)\leq Z_0(\nu)e^{\nu E}\), the norm bound gives
\[
 \frac{|a_n(T)|}{n!}
 \leq Z_0(\nu)\frac{T^n}{2^n(n!)^2}
          \sum_{E\geq0}E^{2n}e^{-tE}
 \leq Z_0(\nu)e^t
       \frac{T^n(2n)!}{2^n(n!)^2t^{2n+1}}.
\]
For the last step, integration over \([E,E+1]\) bounds the
summand by \(e^t\int_E^{E+1}x^{2n}e^{-tx}\,dx\).
The central binomial coefficient has \(n\)th root tending to four.
Consequently
\[
 \limsup_n\left(\frac{|a_n(T)|}{n!}\right)^{1/n}
 \leq\frac{2T}{(T-\nu)^2}.
\]
Letting \(\nu\) decrease to zero gives the upper bound \(2/T\).

For even \(n\), retain one extreme eigenvector at energy
\(E_n=\lfloor2n/T\rfloor\), with \(n\) sufficiently large.
It contributes
\[
 \frac{a_n(T)}{n!}
 \geq\frac{T^n e^{-TE_n}}{(n!)^2}
          \left(\frac{E_n(E_n-1)}2\right)^n.
\]
Stirling's formula makes the \(n\)th root of this expression
tend to \(2/T\). This proves the exact Borel radius.
Multiplication by \(n!\) makes the corresponding coefficient
roots unbounded, proving radius zero for the original series.

For endpoint divergence, retain both extreme vectors at every
energy \(E\geq2\). For \(\zeta>0\), their combined contribution
to the Borel sum is
\[
 e^{-TE}\bigl(I_0(X_E)+J_0(X_E)\bigr),\qquad
 X_E=2\sqrt{T\zeta E(E-1)/2},
\]
where
\[
 I_0(X)=\sum_{n\geq0}\frac{(X^2/4)^n}{(n!)^2},\qquad
 J_0(X)=\sum_{n\geq0}\frac{(-1)^n(X^2/4)^n}{(n!)^2}.
\]
Their sum retains precisely the even powers. At \(\zeta=T/2\),
\(X_E=TE-T/2+O(E^{-1})\). The elementary integral representation
\[
 I_0(X)=\frac1\pi\int_0^\pi e^{X\cos\theta}\,d\theta
\]
gives \(I_0(X)\geq c e^X/\sqrt X\) for large positive \(X\),
by restricting to \(0\leq\theta\leq X^{-1/2}\).
The oscillatory integral for \(J_0\) gives \(|J_0(X)|\leq1\).
Thus the retained contribution is bounded below by a positive
multiple of \(E^{-1/2}\) for large \(E\). Its sum diverges.
All coefficients of the even Borel series are nonnegative, so
this subseries comparison is valid. Evenness gives the other endpoint.
\end{proof}

A fixed vacuum-energy shift multiplies the entire scalar series by
a nonzero constant and changes neither radius. Division by the
free partition function has the same property. The theorem does not
construct continuation outside the disk or choose a Laplace contour.
A positive resummation requires a state functional on an observable
star-algebra, with positivity and sewing for every admitted observable.
Those requirements are stronger than a continuation of the partition
function alone.

\section{Fermionic charges and a local quartic density}

The nonlocal square \(H_0^2\) stabilizes the spectrum, but it does
not express the correction as the integral of a local circle density.
For that purpose, a second realization of the same Fock space exposes
a useful positive charge.

Let \(r\) run over half-integers. A fermionic occupation set differs
from the filled negative sea
\(S_0=\{-1/2,-3/2,\ldots\}\) in finitely many places.
The corresponding wedges form an orthonormal basis. Creation
\(c_r^*\) inserts the mode \(r\), and annihilation \(c_r\)
removes it, with the exterior signs. They satisfy
\(\{c_r,c_s^*\}=\delta_{rs}\).
Normal ordering subtracts the vacuum occupation:
\[
 {:}c_r^*c_s{:}=c_r^*c_s-\delta_{rs}\mathbf1_{r<0}.
\]
The charge is the number of added positive particles minus the
number of holes in the negative sea. Charge-zero occupations are
indexed by partitions \(\mu\), through
\[
 S_\mu=\{\mu_i-i+1/2:i\geq1\}.
\]
Their energy is \(|\mu|\). Write
\(A(\theta)=\sum_{n\ne0}\alpha_n e^{-in\theta}\) for the angular
current, with each Fourier coefficient acting on finite-energy vectors.

\begin{proposition}\label{prop:vol2-boson-fermion-charge}
The charge-zero wedge space is unitarily isomorphic to \(\mathcal F\),
with vacuum sent to the filled sea and
\begin{equation}\label{eq:vol2-fermion-current}
 \alpha_n=\sum_{r\in\mathbb Z+1/2}{:}c_{r-n}^*c_r{:}
 \qquad(n\ne0).
\end{equation}
Define diagonal operators on the partition basis by
\begin{equation}\label{eq:vol2-partition-charges}
 Q_j(\mu)=\sum_{i\geq1}
 \left((\mu_i-i+1/2)^j-(-i+1/2)^j\right).
\end{equation}
Then \(Q_1=H_0\), \(Q_2=2V\), and
\begin{equation}\label{eq:vol2-local-quartic-charge}
 Q_3=\frac14\int_0^{2\pi}\frac{d\theta}{2\pi}
 \left({:}A(\theta)^4{:}
       +{:}(\partial_\theta A(\theta))^2{:}
       -\frac12{:}A(\theta)^2{:}\right).
\end{equation}
All these identities hold on finite-energy vectors.
\end{proposition}

\begin{proof}
For a fixed occupation set, only finitely many summands of
\eqref{eq:vol2-fermion-current} act. The anticommutation relations
give \([\alpha_m,\alpha_n]=m\delta_{m+n,0}\): off-diagonal
matrix units cancel, and for \(m=-n>0\) the vacuum boundary
contributes the \(m\) modes crossing zero. The positive modes
kill the filled sea. The vectors obtained by applying products of
negative modes therefore have the same orthogonal norms
\(z_\mu\) as the bosonic occupation monomials. They are linearly
independent. In energy \(E\), both spaces have dimension \(p(E)\),
so they span. This proves the isometry and its extension to the
Hilbert completions.

The diagonal charges and their vacuum scalar can be calculated
through a generating field. Put \(\varsigma(s)=2\sinh(s/2)\) and
\[
 \mathcal E_n(s)=\sum_r e^{s(r-n/2)}{:}c_{r-n}^*c_r{:}
                   +\frac{\delta_{n0}}{\varsigma(s)}.
\]
Direct multiplication of the fermionic matrix units gives
\begin{equation}\label{eq:vol2-bilocal-commutator}
 [\alpha_m,\mathcal E_n(s)]
       =\varsigma(ms)\mathcal E_{m+n}(s).
\end{equation}
When \(m+n=0\), the boundary sum is
\(\varsigma(ms)/\varsigma(s)\), which accounts for the scalar
in this equation. These are the charge and operator conventions
of Okounkov--Pandharipande \cite{OkounkovPandharipande2002}*{Sections 2.1--2.2}.
Here the vacuum scalar is subtracted from every \(Q_j\).

The normally ordered exponential
\[
 \mathcal B(z,s)=\frac1{\varsigma(s)}
 {:}\exp\!\left(\sum_{n\ne0}
       \frac{\varsigma(ns)}n\alpha_nz^{-n}\right){:}
\]
has commutator
\([\alpha_m,\mathcal B]=\varsigma(ms)z^m\mathcal B\).
The field \(\sum_n\mathcal E_n(s)z^{-n}\) has the same
commutators by \eqref{eq:vol2-bilocal-commutator}.
Both have vacuum matrix element \(1/\varsigma(s)\).
Commuting positive modes from the left and negative modes from the
right reduces any matrix element between created vacuum vectors
to that scalar matrix element. Induction on the total number of
external modes therefore proves equality of the two fields.
This is an identity of formal Laurent series in \(s\), evaluated
between finite-energy vectors.

Let \(\operatorname{CT}_z\) denote the angular constant term.
Taking it in the field identity and subtracting the vacuum gives
\begin{equation}\label{eq:vol2-charge-bilocal}
 \sum_{j\geq0}\frac{s^j}{j!}Q_j
 =\frac1{\varsigma(s)}
 \left[\operatorname{CT}_z
 {:}\exp\!\left(\sum_{n\ne0}
       \frac{\varsigma(ns)}n\alpha_nz^{-n}\right){:}-1\right].
\end{equation}
The left side follows directly from the finite occupation difference
in \eqref{eq:vol2-partition-charges}. In charge zero, \(Q_0=0\).

Write \(A(z)=\sum_{n\ne0}\alpha_nz^{-n}\) and
\(B(z)=\sum_{n\ne0}n^2\alpha_nz^{-n}\). The exponent is
\(sA+s^3B/24+O(s^5)\), while
\(\varsigma(s)^{-1}=s^{-1}-s/24+O(s^3)\).
Since \(\operatorname{CT}A=0\), extraction gives
\begin{align*}
 Q_1&=\tfrac12\operatorname{CT}{:}A^2{:}=H_0,\\
 Q_2&=\tfrac13\operatorname{CT}{:}A^3{:}=2V,\\
 Q_3&=\tfrac14\operatorname{CT}{:}A^4{:}
      +\tfrac14\operatorname{CT}{:}AB{:}
      -\tfrac18\operatorname{CT}{:}A^2{:}.
\end{align*}
Finally,
\[
 \operatorname{CT}{:}AB{:}
 =2\sum_{n>0}n^2\alpha_{-n}\alpha_n
 =\operatorname{CT}{:}(\partial_\theta A)^2{:}.
\]
This proves \eqref{eq:vol2-local-quartic-charge}, including its
derivative correction and quadratic subtraction.
\end{proof}

The isometry in the proof sends the orthonormal Schur basis
\[
 s_\mu(p)=\sum_{\nu\vdash |\mu|}
                  \frac{\chi_\mu(\nu)}{z_\nu}p_\nu
\]
to \(S_\mu\). One can check this directly on finite alternants:
multiplication by \(p_n\) raises an occupied exponent by \(n\),
and reordering supplies the exterior sign. Character orthogonality
gives the stated normalization. Formula~\eqref{eq:vol2-partition-charges}
also gives \(V s_\mu=(\sum_{(i,j)\in\mu}(j-i))s_\mu\),
consistent with the two extreme partitions above.

\section{A positive local preparation}

Let \(d\) be the number of diagonal boxes of \(\mu\), and put
\[
 p_i=\mu_i-i+\tfrac12,\qquad
 h_i=\mu_i'-i+\tfrac12\qquad(1\leq i\leq d).
\]
These positive half-integers are the added particle energies and
the absolute energies of the holes. Separating the two occupation
differences in \eqref{eq:vol2-partition-charges} gives
\begin{equation}\label{eq:vol2-particle-hole-charges}
 E=\sum_i(p_i+h_i),\qquad
 Q_2=\sum_i(p_i^2-h_i^2),\qquad
 Q_3=\sum_i(p_i^3+h_i^3).
\end{equation}
In particular \(Q_3\geq0\), and Cauchy--Schwarz yields
\begin{equation}\label{eq:vol2-charge-inequality}
 Q_2^2\leq E Q_3.
\end{equation}
The two vectors used here have entries
\((\sqrt{p_i},\sqrt{h_i})\) and
\((p_i^{3/2},-h_i^{3/2})\).

\begin{theorem}\label{thm:vol2-local-stable-state}
For real \(g\), define
\begin{equation}\label{eq:vol2-local-stable-hamiltonian}
 H_g=H_0+gV+\frac{g^2}{8}Q_3.
\end{equation}
It is essentially self-adjoint on finite-energy vectors, with its
real diagonal closure in the partition basis, and
\(H_g\geq H_0/2\). For every \(T>0\),
\[
 1\leq Z(T,g)=\operatorname{Tr}e^{-TH_g}\leq Z_0(T/2)<\infty.
\]
The normalized trace is a positive state on bounded operators.
Finite Fourier Wick polynomials also have finite expectations.
The unnormalized heat kernels sew exactly, and their dependence
on real \(g\) is smooth in trace norm.
\end{theorem}

\begin{proof}
All charges are real diagonal in an orthonormal basis. Their sum
has maximal domain
\[
 \left\{(v_\mu):\sum_\mu|H_g(\mu)|^2|v_\mu|^2<\infty\right\}.
\]
Finite-support vectors are a core for this multiplication operator.
By \eqref{eq:vol2-charge-inequality},
\[
 |gV|\leq\frac{|g|}{2}\sqrt{E Q_3}
       \leq\frac E2+\frac{g^2Q_3}{8}.
\]
The last inequality is the nonnegativity of
\((\sqrt{E/2}-|g|\sqrt{Q_3/8})^2\).
It proves the lower bound and hence trace class by summing
\(e^{-TH_g(\mu)}\leq e^{-T|\mu|/2}\).
The vacuum contributes one. Polynomial energy bounds extend the
trace domain to the stated unbounded insertions.

The spectral calculus gives
\(e^{-T_1H_g}e^{-T_2H_g}=e^{-(T_1+T_2)H_g}\).
On compact real coupling intervals, each derivative of order \(k\)
of a diagonal exponential is bounded by a constant times
\((1+E)^{3k}e^{-TE/2}\). Indeed
\(|Q_2|\leq E^2\) and \(Q_3\leq E^3\) follow from
\eqref{eq:vol2-particle-hole-charges}. This dominating sequence is
summable against \(p(E)\). Termwise differentiation is therefore
valid in trace norm, at every order.
\end{proof}

The correction is the integral of the local density
\eqref{eq:vol2-local-quartic-charge}. It first appears at order
\(g^2\). The derivatives of the resulting state agree with the
formal expansion of this modified Hamiltonian.

The real smoothness does not imply a holomorphic trace-class family
near \(g=0\). For \(g=i\gamma\ne0\), the real part of
\(H_g\) is \(H_0-\gamma^2Q_3/8\). On a one-row partition
of size \(E\),
\(Q_3=(E-1/2)^3+1/8\), so the corresponding exponential is
unbounded. Any holomorphic trace-class extension agreeing on a real
interval would have these diagonal matrix entries, by uniqueness
of scalar holomorphic continuation. Such an extension is impossible.

There is also a convergent product expression for the positive trace.
For \(r\in\mathbb Z_{\geq0}+1/2\), set
\[
 \varepsilon_+(r)=r+\frac g2r^2+\frac{g^2}{8}r^3,
 \qquad
 \varepsilon_-(r)=r-\frac g2r^2+\frac{g^2}{8}r^3.
\]
Each is at least \(r/2\). The grand-canonical product
\begin{equation}\label{eq:vol2-local-grand-product}
 \mathcal Z_{T,g}(y)=\prod_{r>0}
 \bigl(1+y e^{-T\varepsilon_+(r)}\bigr)
 \bigl(1+y^{-1}e^{-T\varepsilon_-(r)}\bigr)
\end{equation}
converges absolutely and uniformly on \(|y|=1\). The sum of
the magnitudes of its nonconstant factors is bounded by
\(2\sum_{r>0}e^{-Tr/2}\). Consequently
\begin{equation}\label{eq:vol2-charge-projected-trace}
 Z(T,g)=[y^0]\mathcal Z_{T,g}(y)
       =\frac1{2\pi}\int_0^{2\pi}\mathcal Z_{T,g}(e^{i\theta})\,d\theta.
\end{equation}
Each factor records one particle or one hole. The exponent of \(y\)
records their charge difference, so its zero coefficient retains
exactly the partition basis of \(\mathcal F\). Termwise extraction
is justified by absolute convergence and gives the spectral trace.
At \(g=0\), the expression reduces to \(Z_0(T)\). Omitting the
charge projection would change the Hilbert space.

\section{The sharp coefficient and the critical state}

To determine what the coefficient \(1/8\) achieves, consider
the larger real family
\[
 H_{g,c}=H_0+\frac g2Q_2+c g^2Q_3.
\]
It is defined by the same diagonal closure and local density.

\begin{proposition}\label{prop:vol2-sharp-stabilization}
For \(c>0\), the largest constant \(a_c\) such that
\(H_{g,c}\geq a_cH_0\) for every real \(g\) is
\[
 a_c=1-\frac1{16c}.
\]
Nonnegativity for every coupling holds exactly for \(c\geq1/16\).
The uniform half-energy bound holds exactly for \(c\geq1/8\).
For \(c>1/16\), the heat operators and normalized densities
are continuous in \(g\) in trace norm.
\end{proposition}

\begin{proof}
For a nonempty partition, complete the square:
\[
 H_{g,c}
 =cQ_3\left(g+\frac{Q_2}{4cQ_3}\right)^2
       +E-\frac{Q_2^2}{16cQ_3}
 \geq\left(1-\frac1{16c}\right)E.
\]
The vacuum charges vanish. On the row partition \((N)\), put
\(p=N-1/2\). Then
\[
 E=N,\qquad Q_2=p^2-\tfrac14=N(N-1),\qquad
 Q_3=p^3+\tfrac18.
\]
The ratio \(Q_2^2/(EQ_3)\) tends to one. At
\(g=-Q_2/(4cQ_3)\), the ratio \(H_{g,c}/E\) therefore tends
to \(a_c\), proving optimality. For \(c<0\), the energy of
a nonempty partition becomes negative at large \(|g|\).
For \(c=0\), a row of length at least two has a nonzero cubic
charge and becomes negative for one sign of \(g\).
This proves both coefficient thresholds.

If \(c>1/16\), the diagonal heat entries are bounded uniformly
by \(e^{-Ta_cE}\). Their sum is \(Z_0(Ta_c)\).
For \(g_j\to g\), differences converge entrywise and are
dominated by twice this summable sequence. Their sum of absolute
values is precisely the trace norm. The traces also converge and
are at least one, so division gives continuity of the densities.
\end{proof}

At the critical value, fixed-coupling trace class survives but
this continuity fails.

\begin{theorem}\label{thm:vol2-critical-state-escape}
Let \(c=1/16\), \(K_g(T)=e^{-TH_{g,c}}\), and
\(\rho_g(T)=K_g(T)/\operatorname{Tr}K_g(T)\).
For every real \(g\) and \(T>0\), the heat operator is positive
and trace class. As \(g\to0\), it converges strongly to
\(e^{-TH_0}\). Nevertheless, along
\[
 g_N=-\frac4{N-1/2},\qquad N\geq4,
\]
one has
\begin{align*}
 \operatorname{Tr}K_{g_N}(T)&\longrightarrow\infty,\\
 \|K_{g_N}(T)-K_0(T)\|_1&\longrightarrow\infty,\\
 \|\rho_{g_N}(T)-\rho_0(T)\|_1&\longrightarrow2.
\end{align*}
\end{theorem}

\begin{proof}
The particle and hole contributions are now
\[
 \varepsilon_\pm(r)=r\left(1\pm\frac{gr}{4}\right)^2\geq0.
\]
For fixed \(g\ne0\) and \(r\geq8/|g|\), both are at least
\(r\). Each lower particle or hole mode can occur only once.
Thus \(H_{g,c}\geq E-C_g\), where the finite constant is
\[
 C_g=\sum_{0<r<8/|g|}
 \bigl((r-\varepsilon_+(r))_+
       +(r-\varepsilon_-(r))_+\bigr),\qquad x_+=\max(x,0).
\]
It follows that \(\operatorname{Tr}K_g(T)\leq e^{TC_g}Z_0(T)\).
At \(g=0\), trace class holds directly. Nonnegativity bounds
\(\|K_g(T)\|\) by one. Its entries converge on every fixed
partition vector, so density of their finite span proves strong
convergence on the Hilbert space.

Set \(p=N-1/2\). For the row partitions \((N+k)\), with
\(0\leq k\leq\lfloor\sqrt p\rfloor\), the energy at \(g_N\) is
\begin{equation}\label{eq:vol2-critical-row-energy}
 H_{g_N,1/16}(N+k)
 =\frac{(p+k)k^2}{p^2}
       +\frac12\left(1+\frac1{2p}\right)^2<3.
\end{equation}
The first term is at most \(1+1/\sqrt p<2\), and the second
is less than one for \(p\geq7/2\). These distinct states give
\[
 \operatorname{Tr}K_{g_N}(T)
 \geq(1+\lfloor\sqrt p\rfloor)e^{-3T}\longrightarrow\infty.
\]
The trace difference is a lower bound on the trace-norm difference,
which proves the second limit.

For a fixed finite energy projection \(F\),
\(\operatorname{Tr}(F\rho_{g_N}(T))\to0\), since its numerator
is bounded by \(\operatorname{rank}F\) and the denominator diverges.
Given \(\eta>0\), choose \(F\) with
\(\operatorname{Tr}(F\rho_0(T))>1-\eta\).
Both densities are diagonal and have trace one. Summing absolute
differences on \(F\) and its complement gives
\[
 \|\rho_{g_N}(T)-\rho_0(T)\|_1
 \geq2\left|\operatorname{Tr}(F\rho_{g_N}(T))
              -\operatorname{Tr}(F\rho_0(T))\right|.
\]
The lower limit is at least \(2(1-\eta)\), and the norm is at
most two. Let \(\eta\) decrease to zero.
\end{proof}

The critical example distinguishes strong convergence of heat
operators from convergence of the prepared states. Above the
critical coefficient, uniform energy decay prevents the mass
from escaping to partitions of increasing size. The coefficient
\(1/8\) gives the particularly simple bound \(H_g\geq H_0/2\).

\section{The interacting current and the remaining state map}

Return to the local Hamiltonian \eqref{eq:vol2-local-stable-hamiltonian}.
Its state preserves energy, so a two-current expectation vanishes
unless the two indices sum to zero. The fermionic occupations give
its remaining coefficients explicitly. For \(r>0\), put
\[
 f_r(y)=\frac{y e^{-T\varepsilon_+(r)}}
                   {1+y e^{-T\varepsilon_+(r)}},\qquad
 f_{-r}(y)=\frac1{1+y^{-1}e^{-T\varepsilon_-(r)}}.
\]
These functions record occupancy of a positive mode and occupancy
of a negative mode in the filled sea. Their denominators do not
vanish on \(|y|=1\). Define, for \(n>0\),
\begin{equation}\label{eq:vol2-interacting-current-coefficient}
 A_n(T,g)=
 \frac{[y^0]\left(\mathcal Z_{T,g}(y)
       \sum_{r\in\mathbb Z+1/2}f_{r-n}(y)(1-f_r(y))\right)}{Z(T,g)}.
\end{equation}
For fixed \(n\), the tails decay exponentially and uniformly on
the unit circle. The finitely many modes crossing the sea boundary
are harmless.

\begin{proposition}\label{prop:vol2-interacting-current-pole}
The coefficient \(A_n(T,g)\) is
\(\varphi_{T,g}(\alpha_n\alpha_{-n})\). If
\(c_n=A_n-n\), then \(c_n\geq0\) and
\begin{equation}\label{eq:vol2-interacting-current-bound}
 c_n\leq C_T e^{-Tn/4},\qquad
 C_T=\sum_{E\geq0}E p(E)e^{-TE/4}<\infty.
\end{equation}
The current two-point function is
\begin{equation}\label{eq:vol2-stable-current-kernel}
 \varphi_{T,g}(J(z)J(w))
 =\frac{\lambda}{(z-w)^2}
   +\frac{\lambda}{zw}\sum_{n\geq1}c_n(T,g)
          \bigl((w/z)^n+(z/w)^n\bigr).
\end{equation}
Its radially ordered domain includes \(e^{-T/4}<|w/z|<1\).
The remainder and all fixed holomorphic derivatives extend across
\(z=w\ne0\).
\end{proposition}

\begin{proof}
In the grand-canonical occupation measure, a raising move from
\(r-n\) to \(r\) has weight \(f_{r-n}(1-f_r)\).
By \eqref{eq:vol2-fermion-current}, the squared norm of this move
is the corresponding contribution to \(\alpha_n\alpha_{-n}\).
Different moves from a fixed occupation vector give orthogonal
vectors. Summing their expectations therefore gives the numerator
in \eqref{eq:vol2-interacting-current-coefficient}. The insertion
is charge neutral, so taking the zero-charge coefficient commutes
with it. Absolute convergence justifies the extraction.

The commutator gives
\(c_n=\varphi_{T,g}(\alpha_{-n}\alpha_n)\geq0\).
This positive operator is at most \(H_0\), and it kills energy
spaces below \(n\). The half-energy bound and \(Z(T,g)\geq1\)
give
\[
 c_n\leq\sum_{E\geq n}E p(E)e^{-TE/2}
       \leq e^{-Tn/4}\sum_{E\geq0}E p(E)e^{-TE/4}.
\]
The constant is finite by the free partition estimate.
Inserting \(A_n=n+c_n\) into the current series yields
\eqref{eq:vol2-stable-current-kernel}. The vacuum sum supplies
the pole; the exponential bound makes the two remainder series
normally convergent in a neighborhood of \(|w/z|=1\).
\end{proof}

The modified preparation retains the current pole while changing
the regular coefficients. It need not be Gaussian as a state of
the bosonic current algebra. Its operator sewing and scalar
normalization are fixed by the spectral preparation, including
the charge projection. A conformal shift \(-1/24\) multiplies
its unnormalized kernels by \(e^{T/24}\) and leaves normalized
correlations unchanged.

For the unchanged cubic interaction, the outstanding construction
is a state-valued Borel continuation with a specified Laplace
prescription on a chosen observable star-algebra and sewing
completion. It must give a normalized positive functional on every
allowed observable and preserve the unnormalized cylinder
composition. The scalar Borel boundary does not decide whether
such a prescription exists. The positive states of the modified
local Hamiltonian differ from that interaction at order \(g^2\).

A bulk realization has a further requirement. It must construct
the boundary preparation from independently defined bulk fields,
counterterms and a boundary condition, then intertwine the current,
collision and sewing operations on their declared domains.
An equality of characters or a matching determinant alone does
not supply that state map.
